\section{Further extensions worth trying}\label{sec:further}

These were not fitted here. Each stays close to the current models and uses data that already exist.

\begin{enumerate}[leftmargin=1.6em]
\item \textbf{Use the frequency-probe trials.}
Vera's selection drops 3,232 completed trials in which a single vibration was given at 150 or 250\,Hz instead of 200\,Hz (all unimanual, from the same days).
With a frequency term in $\eta$ they would add about 75 unimanual trials per session, which would sharpen each session's bias, slope and unimanual lapses: the baseline against which every bimanual effect is measured.
They also test whether Sirius judges amplitude or perceived intensity, which bears on whether the distractor adds to the target's intensity.

\item \textbf{Response times.}
The trial files record when the choice targets appear and when the response movement starts (median response time 562\,ms).
A first look supports treating some bimanual ``lapses'' as real decisions driven by the distractor (Section~\ref{sec:interp}).
A drift-diffusion version of Model~E, in which the distractor adds to the drift rate and lapses are a separate response process, would use both choices and times and could tell a change in the evidence from a change in the criterion.

\item \textbf{A more flexible trend in the criterion.}
The linear trend and the amplitude term fit equally well. A random walk over sessions would let the criterion wander without assuming a shape, and a criterion that follows a running average of recent distractor amplitudes would describe slow adaptation. On these data they would be hard to tell apart for the same reason as D+crit and D+order, but they are the right models for data in which amplitude and time are not confounded.

\item \textbf{Trial history.}
A logistic check with the full session-by-hand structure finds only a weak tendency to switch away from the previous choice (logit $-0.20\pm0.11$) and no clear win-stay effect ($0.13\pm0.12$).
Adding the previous choice to $\eta$ is cheap, but it is unlikely to change the conclusions.

\item \textbf{Correlated session effects.}
The models draw each session parameter independently. A full covariance across a session's parameters (an LKJ prior on bias, slope and lapses, and across the two hands) is the general form of the shared daily bias tested in D+shared.

\item \textbf{More animals.}
A monkey level above the session level would separate effects that hold for every animal from the individual pattern seen here, such as the right hand's extra shift.
\end{enumerate}
